Results are shown in Fig. \ref{fig:ExtPin_group}. The initial identification was fitted in cross-validated rounds over the \qty{48.0}{\cm}, \qty{45.7}{\cm}, and \qty{41.5}{\cm} actuators, holding each actuator out in turn, and the quoted round is the one holding out $l_{0}=\qty{41.5}{\cm}$; it resulted in $\chi_{0} = \qty{-10.9}{\mm}$ and $\chi_{3} = 0.123$, the extensor length offset $\chi_{0}$ being nearly equal in magnitude to the flexor offset but opposite in sign, and the goodness-of-fit values of that initial round are found in Table~\ref{tab:TorqueFit}. The extensor identification of record supersedes it: the two-stage refit of Section~\ref{sec:identification} locks the stiffness pair at the flexor-identified values and solves $\chi_{0} = \qty{-6.38}{\mm}$ and $\chi_{3} = 0.158$ (Table~\ref{tab:AdoptedXi}) on the six-test pool (Table~\ref{tab:ExtPinGof}) drawn from the same test campaign, with the \qty{43.6}{\cm}, \qty{43.2}{\cm}-tendon, and \qty{45.7}{\cm} tests in training and the \qty{40.0}{\cm} plain, \qty{40.0}{\cm}-tendon, and \qty{46.5}{\cm} tests held out for validation; Fig.~\ref{fig:ExtPin_group} shows the six pool tests at the adopted terms.
